\chapter{Nash Equilibrium}

\section{Nash Equilibrium and Nash's Theorem}

\defn{Nash Equilibrium}{
  A strategy profile $\pr{\bar\sigma_1,\bar\sigma_2,\cdots,\bar\sigma_n}$ (where $\bar\sigma_i\in \Delta(S_i)$) is a Nash equilibrium of the game $G = \pr{S_i,u_i}_{i = 1}^n$ if for any player $i$,
  \[
  u_i(\bar\sigma_i,\bar\sigma_{-i})\ge u_i(s_i,\bar\sigma_{-i}),\quad \forall s_i\in S_i.
  \]
}

\rmkb[]{
  Equivalently we can define Nash equilibrium by
  \[
  u_i(\bar\sigma_i,\bar\sigma_{-i})\ge u_i(\sigma_i,\bar\sigma_{-i}),\quad \forall \sigma_i\in \Delta(S_i).
  \]
  But it suffices to only check pure strategies $s_i$ because mixed strategies are just probability distributions over pure strategies, and the expected utility of a mixed strategy is a convex combination of the utilities of the pure strategies.
}

\thmp{Nash}{
  Every finite game has a Nash equilibrium in mixed strategies.
}{
  Here we apply the Kakutani's Fixed Point Theorem.
  \thm{Kakutani's Fixed Point Theorem}{
    Suppose $F:X\rightrightarrows X$ (i.e., $F(x)\subseteq X$), and 
    \begin{itemize}
      \item $X$ is compact and convex;
      \item $F(x)$ is convex for all $x\in X$;
      \item $F$ has a closed graph, i.e., the set $\brace{(x,y)|f\in F(x)}$ is closed.
    \end{itemize}
    Then, there exists $x^*\in X$ such that $x^*\in F(x^*)$.
  }
  Let $X_i\equiv \Delta(S_i)$, and $X = X_1\times X_2\times \cdots \times X_n = \prod_{i = 1}^{n}\Delta(S_i)$. Define the best response correspondence $B_i:\prod_{j\ne i}\Delta(S_j) \rightrightarrows \Delta(S_i)$ by
  \[
  \sigma_i\in B_i(\sigma_{-i}) \iff u_i(\sigma_i,\sigma_{-i})\ge u_i(\sigma_i',\sigma_{-i}),\quad \forall \sigma_i'\in \Delta(S_i).
  \]
  And we define $B:X\rightrightarrows X$ by 
  \[
  B(\sigma) = \begin{pmatrix}
    B_1\\
    B_2\\
    \vdots\\
    B_n
  \end{pmatrix}(\sigma).
  \]
  It is easy to verify that $B(\sigma)$ is convex for all $\sigma\in X$, and $B$ has a closed graph. Therefore, by Kakutani's theorem, there exists $\bar\sigma\in X$ such that $\bar\sigma\in B(\bar\sigma)$, which means that $\bar\sigma$ is a Nash equilibrium.
}
\rmkb[]{
  \begin{itemize}
    \item ``Mixed strategies'' is important for Nash's Theorem. Finite games may not have NE in pure strategies.
    \item If $\bar\sigma$ is a NE, and $\bar\sigma_i(s_i)>0$, then $s_i$ must be a best response to $\bar\sigma_{-i}$, and thus $u_i(s_i,\bar\sigma_{-i})=u_i(\bar\sigma_i,\bar\sigma_{-i})$. In other words, if a pure strategy is played with positive probability in a mixed strategy NE (i.e., in the \textit{support} of the mixed strategy $\bar\sigma$), then it must yield the same expected payoff as the mixed strategy itself.
    \item Then why bother mixing? Because each player mixes to make other players mix. If a player chooses a pure strategy, then other players will have no incentive to mix, and thus the first player will have no incentive to mix either. Therefore, mixing can be a strategic move to induce other players to mix, which can lead to a more favorable outcome for the player.
  \end{itemize}
}

\section{Potential Games}

One of the central questions in game theory is whether a game has a \textit{pure strategy Nash equilibrium} (PSNE). Some games, like matching penny, have no PSNE. However, certain classes of games are guaranteed to have PSNE. Two important classes are:
\begin{itemize}
  \item Potential games
  \item Supermodular games
\end{itemize}

In this section we will focus on potential games, which are a broad class of games that include many interesting examples, such as congestion games, coordination games, and certain types of auctions. Potential games have a special structure that allows us to analyze them using a potential function, which captures the incentives of all players in a single function.

\subsection{Definition}

\defn{Potential Game}{
  A finite game $G = (S_i, u_i)_{i = 1}^n$ is called a \textit{potential game} if there exists a function $P: S \to \RR$ (where $S = S_1 \times S_2 \times \cdots \times S_n$) such that for every player $i$ and every pair of strategies $s_i, s_i' \in S_i$:
  \[
    u_i(s_i, s_{-i}) - u_i(s_i', s_{-i}) = P(s_i, s_{-i}) - P(s_i', s_{-i}),\quad s_{-i} \in S_{-i}.
  \]
  The function $P$ is called the \textit{potential function}.
}

\ex[Congestion Game]{
  Consider a road network where $n$ drivers choose routes. Let $s_i$ be the route chosen by driver $i$. The payoff to driver $i$ is:
  \[
    u_i(s_i, s_{-i}) = -c(s_i, n_i)
  \]
  where $c(r, n_r)$ is the congestion cost on route $r$ when $n_r$ drivers use it. Define the potential function as:
  \[
    P(s) = -\sum_{r} \int_0^{n_r(s)} c(r, x) dx
  \]
  This is indeed a potential function because the private incentive to switch routes equals the change in total cost (up to a sign).
}

\rmkb{
  \begin{itemize}
    \item Potential games can be thought of as games where all players have identical payoffs.
    \item Potential games arise naturally when:
\begin{itemize}
  \item Players care about a common objective (like aggregate welfare)
  \item Individual incentives are aligned with some global measure
  \item Strategic interactions are ``symmetric'' in the sense of cross-partial derivatives
\end{itemize}
    \item The potential function $P$ captures the payoff differences for individual players when they deviate from one strategy to another. The key insight is that a player's incentive to deviate depends only on how the potential function changes, not on the baseline payoff levels.
    \item In the continuous case with smooth payoff functions, a game is a potential game if and only if:

\[
  \begin{aligned}
    \frac{\partial P}{\partial s_i} = \frac{\partial u_i}{\partial s_i} \quad \text{for all players } i
  \end{aligned}
\]

This leads to the following necessary condition, for all pairs of players $i$ and $j$:
\[
  \frac{\partial^2 u_i}{\partial s_j \partial s_k} = \frac{\partial^2 u_j}{\partial s_k \partial s_j} \quad \text{whenever these derivatives exist.}
\]
  \end{itemize}
}

\subsection{Existence of PSNE in Potential Games}

\thmp{Existence of PSNE in Potential Games}{
  Every potential game has at least one pure strategy Nash equilibrium.
}{
  A PSNE of game $G$ is a maximizer of the potential function $P$. Since $S$ is finite (in the discrete case), $P$ attains its maximum. At such a maximum point $s^*$, if player $i$ deviates to any other strategy $s_i \neq s_i^*$, then:
\[
  u_i(s_i^*, s_{-i}^*) - u_i(s_i, s_{-i}^*) = P(s_i^*, s_{-i}^*) - P(s_i, s_{-i}^*) \geq 0
\]
Thus no player wants to deviate, so $s^*$ is a PSNE.
}

\rmkb[]{
Different information structures (simultaneous move vs. sequential move) can lead to different equilibria in general games. However, in potential games, the set of pure strategy Nash equilibria is often robust across different information structures. This is because the potential function captures the incentives of all players in a way that is independent of the timing of moves. As a result, the same strategy profiles that maximize the potential function will be equilibria regardless of whether players move simultaneously or sequentially.
}

\section{Supermodular Games}

While potential games ensure the existence of PSNE through a global potential function, supermodular games approach the problem differently: they rely on the structure of strategic complementarities between players' actions.

The key insight is \textit{strategic complementarity}: if one player's action increases, it becomes more attractive for other players to increase their actions as well. This creates a natural coordination mechanism that leads to equilibrium. Unlike potential games, which require a global objective function, supermodular games only require local complementarity conditions on payoffs.

Examples of strategic complementarity abound:
\begin{itemize}
  \item In a price competition setting, if a competitor raises prices, your prices become more attractive to consumers, so you want to raise your price too.
  \item In technology adoption, if more people adopt a technology, the benefits of adoption increase for you (network effects).
  \item In public goods games, if others contribute more, your optimal contribution may also increase.
\end{itemize}

\subsection{Definitions}

We need lattice-theoretic concepts to formalize strategic complementarity. 

\defn{Meet \& Join}{
  Let $x$ and $y$ be two vectors of $\RR^n$. The \textit{meet} of $x$ and $y$, denoted by $x\meet y$, is the coordinate-wise minimum:
  \[
    x\meet y = \begin{pmatrix}
      \min(x_1,y_1)\\
      \min(x_2,y_2)\\
      \vdots\\
      \min(x_n,y_n)\\
    \end{pmatrix}
    \]
    The \textit{join} of $x$ and $y$, denoted by $x\join y$, is the coordinate-wise maximum:
    \[
    x\join y = \begin{pmatrix}
      \max(x_1,y_1)\\
      \max(x_2,y_2)\\
      \vdots\\
      \max(x_n,y_n)\\
    \end{pmatrix}
    \]
}

\rmkb{
  Meet and join decompose a pair of strategy profiles into their ``lower bound'' and ``upper bound'' in each dimension. Later we will see, if a payoff function is supermodular, then being at both extremes (the largest and smallest actions) together is at least as good as being at the two middle ground positions.
}

\defn{Supermodular Function}{
  A function $f: \RR^n \to \RR$ is \textit{supermodular} if for all $x, y \in \RR^n$:
  \[
    f(x\meet y) + f(x\join y) \geq f(x) + f(y)
  \]
  This is called the \textit{supermodularity inequality}.
}

Intuitively, the supermodularity inequality says: having some coordinates large and others small is better than having all coordinates at intermediate levels.

In a game setting, we say player $i$'s payoff is supermodular if:
\[
  u_i(s) + u_i(s') \leq u_i(s \meet s') + u_i(s \join s') \quad \forall s, s' \in S
\]

This means: player $i$'s payoff is supermodular in their own strategy relative to opponents' strategies.

\rmkb[]{
  When $f$ is twice continuously differentiable, supermodularity is equivalent to:

\[
  \frac{\partial^2 f}{\partial x_i \partial x_j} \geq 0 \quad \text{for all } i \neq j.
\]
This means: if you increase $x_i$, the marginal benefit of increasing $x_j$ also increases (\textit{strategic complementarity}).
}



\subsection{Existence of PSNE in Supermodular Games}

The next two lemmas establish that best response correspondences have monotonicity properties under supermodularity, which enables us to use fixed point theorems.

\lemp{Monotonicity of Best Responses}{
  Suppose $u:[0,1]^p \times [0,1]^q \to \RR$ is supermodular. Suppose $y', y'' \in [0,1]^q$ with $y'' \geq y'$. Let $x' \in [0,1]^p$ be a maximizer of $u(\cdot, y')$, and $x'' \in [0,1]^p$ be a maximizer of $u(\cdot, y'')$. Then $x' \meet x''$ is a maximizer of $u(\cdot, y')$, and $x' \join x''$ is a maximizer of $u(\cdot, y'')$.
}{
  By construction, since $x'$ maximizes $u(\cdot, y')$ and $x''$ maximizes $u(\cdot, y'')$:
  \[
    u(x', y') \geq u(x' \meet x'', y')
  \]
  
  Also note that $u(x' \meet x'', y') = u(x' \meet x'', y' \meet y'')$ since $y' \meet y'' = y'$ (as $y' \leq y''$).
  
  Similarly:
  \[
    u(x'', y'') \geq u(x' \join x'', y'') = u(x' \join x'', y' \join y'')
  \]
  
  By supermodularity:
  \[
    u(x', y') + u(x'', y'') \geq u(x' \meet x'', y' \meet y'') + u(x' \join x'', y' \join y'')
  \]
  
  Suppose $u(x', y') > u(x' \meet x'', y' \meet y'')$. Then:
  \[
    u(x', y') + u(x'', y'') > u(x' \meet x'', y' \meet y'') + u(x' \join x'', y' \join y'')
  \]
  which violates supermodularity.
  
  Therefore, we must have:
  \[
    \begin{aligned}
      u(x', y') & = u(x' \meet x'', y' \meet y'')\\
    u(x'', y'') & = u(x' \join x'', y' \join y'')
    \end{aligned}
  \]
  
  This means both $x' \meet x''$ and $x' \join x''$ are maximizers at the respective points.
}

\lemp{Monotonicity of Largest Maximizers}{
  Suppose $u:[0,1]^p \times [0,1]^q \to \RR$ is supermodular. Define:
  \[
    h(y) = \max\{x : x \text{ maximizes } u(\cdot, y)\}
  \]
  (the largest maximizer). Then $h$ is non-decreasing: if $y'' \geq y'$, then $h(y'') \geq h(y')$.
}{
  Suppose $y'' \geq y'$. Let $x' = h(y')$ and $x'' = h(y'')$ be the largest maximizers at $y'$ and $y''$ respectively. By the previous lemma, $x' \meet x''$ maximizes $u(\cdot, y')$ and $x' \join x''$ maximizes $u(\cdot, y'')$.
  
  Since $x'\join x''$ is a maximizer at $y''$ and $x''$ is the largest maximizer at $y''$, we have $x'' \geq x' \join x''$. By definition, $x'\join x'' \geq x'$. So $x'' \geq x'$. Therefore, $h(y'') = x'' \geq x' = h(y')$.
}

\thmp{Existence of PSNE in Supermodular Games}{
  Every finite supermodular game has at least one pure strategy Nash equilibrium.
}{
  
This is a direct result of Tarski's Fixed Point Theorem. 

  \thm{Tarski's Fixed Point Theorem}{If $f: L \to L$ is a non-decreasing function on a complete lattice $L$, then $f$ has a fixed point.}

  For each player $i$, define $b_i(s_{-i})$ as the largest maximizer of $u_i(\cdot, s_{-i})$. By the lemmas above, $b_i$ is non-decreasing in $s_{-i}$.
  
  Define the best response mapping $b: S \to S$ by:
  \[
    b(s) = (b_1(s_{-1}), b_2(s_{-2}), \ldots, b_n(s_{-n}))
  \]
  
  Since $b$ is component-wise non-decreasing and $S$ is a finite lattice, by \textit{Tarski's Fixed Point Theorem}, $b$ has a fixed point $\bar{s}$, i.e., $b(\bar{s}) = \bar{s}$. This fixed point is a Nash equilibrium.
}

\ex[Bertrand Duopoly with Differentiated Products]{
  Consider two firms competing in prices. Let $D_i(p_1, p_2)$ be firm $i$'s demand, which is decreasing in its own price $p_i$ and increasing in the competitor's price $p_j$. With constant marginal cost $c$, firm $i$'s profit is:
  \[
    \pi_i(p_1, p_2) = (p_i - c) D_i(p_1, p_2)
  \]
  
  The key observation is:
  \[
    \frac{\partial^2 \pi_i}{\partial p_i \partial p_j} = \frac{\partial}{\partial p_j}\left[\frac{\partial \pi_i}{\partial p_i}\right] > 0
  \]
  
  This is because when the competitor's price increases, the marginal profit from raising your own price increases (since demand becomes more favorable). Thus, firm $i$'s profit is supermodular in $(p_1, p_2)$, and the duopoly game has a PSNE.
}

